% Auto-generated complete chapter from 29B_demo_reels_benchmark.md
\chapter{29B\_demo\_reels\_benchmark}
\label{complete:root_038}

\begin{sourcemeta}
\textbf{Fuente:} \href{file:///E:/Laboral/29B_demo_reels_benchmark.md}{29B\_demo\_reels\_benchmark.md}\\
\textbf{Seccion:} root\\
\textbf{Estado:} canonical\\
\textbf{Rol:} Modulo final complementario\\
\textbf{Lineas originales:} 882\\
\textbf{Ultima modificacion:} 2026-06-11 18:20\\
\textbf{Deduplicacion conservadora:} 0 bloques repetidos omitidos; 0 caracteres repetidos no reimpresos.
\end{sourcemeta}

\section*{Resumen de orientacion}
El archivo anterior \textbf{no estaba completo} para el propósito real de \texttt{29B}.

\section*{Contenido completo depurado}
\addcontentsline{toc}{section}{Contenido completo depurado}




\subsection{0. Estado del módulo}




\begin{mdcode}
Bloque de codigo: text
Module: 29B_demo_reels_benchmark.md
Input task: F1_demo_reels_benchmark
Date generated: 2026-06-11
Status: Complete v1
Source type: public videos, public portfolio pages, public company demo/case pages
Use: TwinSight X500 demo video planning and portfolio presentation
\end{mdcode}




\subsection{1. Veredicto de completitud}




El archivo anterior \textbf{no estaba completo} para el propósito real de \texttt{29B}.




Funcionaba como un complemento general de casos de estudio, pero no como un módulo autónomo de \textbf{demo reels / video benchmark}. Le faltaban:




\begin{itemize}
\item separación clara entre demo reels personales, reels corporativos y casos de estudio con demo;
\item análisis de duración;
\item análisis de apertura;
\item secuencia de shots;
\item uso de captions;
\item uso de voz en off;
\item uso de breakdown técnico;
\item extracción de patrones aplicables al video de TwinSight;
\item shot list recomendado;
\item estructura de edición;
\item guion base;
\item criterios de calidad;
\item errores a evitar.
\end{itemize}




Este archivo reemplaza al 29B anterior.




\medskip\hrule\medskip




\subsection{2. Objetivo del benchmark}




Este módulo analiza cómo se presentan públicamente proyectos comparables o parcialmente comparables a \textbf{TwinSight X500}, especialmente en video.




El objetivo no es copiar contenido visual, sino extraer patrones de presentación para diseñar un demo video profesional de 60–120 segundos para un perfil de:




\begin{mdcode}
Bloque de codigo: text
Real-Time 3D Developer / Unity Technical Artist
Unity WebGL · CAD-to-Realtime · Technical Visualization · Interactive 3D
\end{mdcode}




\medskip\hrule\medskip




\subsection{3. Perfil de referencia: TwinSight X500}




TwinSight X500 debe presentarse como:




\begin{mdcode}
Bloque de codigo: text
A Unity WebGL technical visualization prototype for drone assembly inspection.
\end{mdcode}




La demo debe comunicar:




\begin{itemize}
\item problema de documentación 2D;
\item solución en 3D interactivo;
\item selección de componentes;
\item exploded view;
\item clipping / cross-section;
\item modos visuales;
\item UI técnica;
\item pipeline CAD-to-realtime;
\item optimización;
\item validación con usuarios;
\item propiedad técnica honesta, incluyendo AI-assisted implementation si se menciona.
\end{itemize}




\medskip\hrule\medskip




\subsection{4. Limitaciones de la recolección}




La tarea F1 produjo una base útil, pero no perfecta.




\begin{center}
\scriptsize
\begin{longtable}{>{\raggedright\arraybackslash}p{0.455\linewidth} >{\raggedright\arraybackslash}p{0.455\linewidth}}
\toprule
Limitación & Impacto \\
\midrule
\endfirsthead
\toprule
Limitación & Impacto \\
\midrule
\endhead
YouTube oculta datos sin sesión & URLs exactas o descripciones pueden estar incompletas \\
Vimeo bloqueó acceso & algunos reels no fueron verificables \\
ArtStation bloqueó por seguridad & algunos reels no fueron accesibles \\
Muchos reels son visuales, no técnicos & falta detalle de stack \\
Empresas publican casos, no reels puros & se usan como referencia de video/case structure \\
No todos muestran duración & se marca como not found \\
\bottomrule
\end{longtable}
\end{center}




\subsection{Regla de uso}




Este benchmark debe usarse como guía de edición y estructura, no como evidencia salarial ni de demanda laboral.




\medskip\hrule\medskip




\section{5. Source inventory}




\subsection{5.1 Fuentes principales}




\begin{center}
\scriptsize
\begin{longtable}{>{\raggedright\arraybackslash}p{0.225\linewidth} >{\raggedright\arraybackslash}p{0.225\linewidth} >{\raggedright\arraybackslash}p{0.225\linewidth} >{\raggedright\arraybackslash}p{0.225\linewidth}}
\toprule
\# & Source label & Type & URL \\
\midrule
\endfirsthead
\toprule
\# & Source label & Type & URL \\
\midrule
\endhead
1 & Robin Cai Visual Effects \& Technical Art Demoreel 2025 & Personal reel & [YouTube search](\url{https://www.youtube.com/results?search_query=Robin+Cai+2025+Visual+Effects+\%26+Technical+Art+Demoreel}) \\
2 & Unity Developer Demo Reel 2025 & Personal reel & [YouTube search](\url{https://www.youtube.com/results?search_query=Unity+Developer+Demo+Reel+2025+Multiplayer+WebGL+Technical+Art+Environment+Art}) \\
3 & Technical Visualization Demo Reel — Kinetic Vision & Corporate reel & [YouTube search](\url{https://www.youtube.com/results?search_query=Kinetic+Vision+Technical+Visualization+Demo+Reel}) \\
4 & Kinetic Vision Technical Visualization service & Corporate service page & [Website](\url{https://kinetic-vision.com/technical-visualization/}) \\
5 & Industrial3D recent animations and demo reels & Corporate reel page & [Website](\url{https://industrial3d.com/3d-animation-in-new-york-ny/}) \\
6 & Work Spatial — AI-Driven Onboarding Platform & Portfolio case demo & [Website](\url{https://virtualverse.studio/projects/work-spatial/}) \\
7 & NBK Virtugate — WebGL + VR Banking & Portfolio case demo & [Website](\url{https://virtualverse.studio/projects/nbk-virtugate/}) \\
8 & Empathy Lab — UK Rail VR Training & Portfolio case demo & [Website](\url{https://virtualverse.studio/projects/empathy-lab/}) \\
9 & WebXR Product Configurator & Case study demo & [Website](\url{https://hapzxr.com/case-studies/webxr-product-configurator/}) \\
10 & Digital Twin of Production Facility & Case study demo & [Website](\url{https://hapzxr.com/case-studies/digital-twin-of-production-facility/}) \\
11 & BIM Site Overlay Visualization & Case study demo & [Website](\url{https://hapzxr.com/case-studies/bim-site-overlay-visualization/}) \\
12 & Construction Confined Space Simulation & Case study demo & [Website](\url{https://hapzxr.com/case-studies/construction-confined-space-simulation/}) \\
\bottomrule
\end{longtable}
\end{center}




\subsection{5.2 Source quality}




\begin{center}
\scriptsize
\begin{longtable}{>{\raggedright\arraybackslash}p{0.305\linewidth} >{\raggedright\arraybackslash}p{0.305\linewidth} >{\raggedright\arraybackslash}p{0.305\linewidth}}
\toprule
Source type & Usefulness & Risk \\
\midrule
\endfirsthead
\toprule
Source type & Usefulness & Risk \\
\midrule
\endhead
Personal technical art reel & High & may be visual-heavy \\
Unity/WebGL personal reel & High & may be game-heavy \\
Corporate visualization reel & Medium-high & may be prerendered \\
Case-study demo page & High & not always video-first \\
VR training case & Medium & weaker WebGL fit \\
Digital twin case & High & may imply larger team/product \\
\bottomrule
\end{longtable}
\end{center}




\medskip\hrule\medskip




\section{6. Benchmark: demo examples}




\subsection{6.1 Personal technical reels}




\begin{center}
\scriptsize
\begin{longtable}{>{\raggedright\arraybackslash}p{0.225\linewidth} >{\raggedright\arraybackslash}p{0.225\linewidth} >{\raggedright\arraybackslash}p{0.225\linewidth} >{\raggedright\arraybackslash}p{0.225\linewidth}}
\toprule
Demo & Role target & Duration & Relevance \\
\midrule
\endfirsthead
\toprule
Demo & Role target & Duration & Relevance \\
\midrule
\endhead
Robin Cai 2025 Visual Effects \& Technical Art Demoreel & Technical Artist & 1:42 & Medium \\
Unity Developer Demo Reel 2025 & Unity / WebGL Developer & 1:17 & High \\
\bottomrule
\end{longtable}
\end{center}




\subsubsection{Notes}




Personal reels tend to:




\begin{itemize}
\item start fast;
\item show work immediately;
\item use short clips;
\item rely on captions or title cards;
\item avoid long explanations;
\item show breadth of skill;
\item end with name/contact/portfolio.
\end{itemize}




\subsubsection{Application to TwinSight}




TwinSight should not behave like a general art reel. It should be closer to a \textbf{technical product demo} with reel pacing.




Use:




\begin{mdcode}
Bloque de codigo: text
fast visual proof + captions + metrics + short technical breakdown
\end{mdcode}




Do not use:




\begin{mdcode}
Bloque de codigo: text
music-only montage with unclear functionality
\end{mdcode}




\medskip\hrule\medskip




\subsection{6.2 Corporate technical visualization reels}




\begin{center}
\scriptsize
\begin{longtable}{>{\raggedright\arraybackslash}p{0.225\linewidth} >{\raggedright\arraybackslash}p{0.225\linewidth} >{\raggedright\arraybackslash}p{0.225\linewidth} >{\raggedright\arraybackslash}p{0.225\linewidth}}
\toprule
Demo & Company & Domain & Relevance \\
\midrule
\endfirsthead
\toprule
Demo & Company & Domain & Relevance \\
\midrule
\endhead
Technical Visualization Demo Reel & Kinetic Vision & Product/process/system visualization & High \\
Interactive 3D Media Reel & Industrial3D & Industrial 3D / simulation / product & Medium-high \\
2022 Demo Reel & Industrial3D & Industrial animation & Medium \\
Bechtel Digital Twin video & Industrial3D & Digital twin / engineering & Medium-high \\
\bottomrule
\end{longtable}
\end{center}




\subsubsection{Notes}




Corporate reels tend to:




\begin{itemize}
\item show many projects quickly;
\item emphasize industry credibility;
\item use polished visuals;
\item show machines, cutaways, assemblies, flows, or environments;
\item focus on technical communication;
\item rarely show implementation details.
\end{itemize}




\subsubsection{Application to TwinSight}




TwinSight should borrow:




\begin{itemize}
\item clarity;
\item object close-ups;
\item section/cutaway shots;
\item before/after visuals;
\item captions that explain value.
\end{itemize}




TwinSight should not copy:




\begin{itemize}
\item purely cinematic pacing;
\item photorealism-first framing;
\item lack of technical ownership;
\item generic marketing language.
\end{itemize}




\medskip\hrule\medskip




\subsection{6.3 WebGL / browser-based case demos}




\begin{center}
\scriptsize
\begin{longtable}{>{\raggedright\arraybackslash}p{0.225\linewidth} >{\raggedright\arraybackslash}p{0.225\linewidth} >{\raggedright\arraybackslash}p{0.225\linewidth} >{\raggedright\arraybackslash}p{0.225\linewidth}}
\toprule
Case & Platform & Core value & Relevance \\
\midrule
\endfirsthead
\toprule
Case & Platform & Core value & Relevance \\
\midrule
\endhead
Work Spatial & WebGL & browser-based XR onboarding & High \\
NBK Virtugate & WebGL + VR & browser onboarding + optional VR & High \\
WebXR Product Configurator & WebGL / WebAR & product inspection/configuration & Very high \\
\bottomrule
\end{longtable}
\end{center}




\subsubsection{Notes}




Browser-based demos are the closest reference for TwinSight because they prove:




\begin{itemize}
\item no-install access;
\item interactive 3D in browser;
\item real-time interaction;
\item fast comprehension;
\item user-facing workflows.
\end{itemize}




\subsubsection{Application to TwinSight}




The video should explicitly show:




\begin{mdcode}
Bloque de codigo: text
Runs in browser
Interactive component selection
Real-time camera movement
Technical UI
No-install viewer
\end{mdcode}




\medskip\hrule\medskip




\subsection{6.4 Digital twin / simulation case demos}




\begin{center}
\scriptsize
\begin{longtable}{>{\raggedright\arraybackslash}p{0.225\linewidth} >{\raggedright\arraybackslash}p{0.225\linewidth} >{\raggedright\arraybackslash}p{0.225\linewidth} >{\raggedright\arraybackslash}p{0.225\linewidth}}
\toprule
Case & Domain & Core feature & Relevance \\
\midrule
\endfirsthead
\toprule
Case & Domain & Core feature & Relevance \\
\midrule
\endhead
Digital Twin of Production Facility & industrial operations & live IoT + 3D facility & High \\
BIM Site Overlay Visualization & construction & BIM + overlay & Medium-high \\
Construction Confined Space Simulation & safety training & guided procedural training & Medium \\
\bottomrule
\end{longtable}
\end{center}




\subsubsection{Notes}




These cases are valuable for:




\begin{itemize}
\item explaining spatial context;
\item showing how 3D replaces fragmented documentation;
\item presenting workflows;
\item connecting interaction to user decisions;
\item using results/metrics.
\end{itemize}




\subsubsection{Application to TwinSight}




TwinSight should use similar language:




\begin{mdcode}
Bloque de codigo: text
spatial context
guided inspection
component-level understanding
interactive technical support
reduced cognitive load
\end{mdcode}




But avoid claiming:




\begin{mdcode}
Bloque de codigo: text
live IoT
predictive maintenance
production digital twin
VR training certification
\end{mdcode}




\medskip\hrule\medskip




\section{7. Video pattern extraction}




\subsection{7.1 Ideal duration}




\begin{center}
\scriptsize
\begin{longtable}{>{\raggedright\arraybackslash}p{0.305\linewidth} >{\raggedright\arraybackslash}p{0.305\linewidth} >{\raggedright\arraybackslash}p{0.305\linewidth}}
\toprule
Video type & Common duration & Recommended for TwinSight \\
\midrule
\endfirsthead
\toprule
Video type & Common duration & Recommended for TwinSight \\
\midrule
\endhead
Personal technical reel & 60–120 s & 90 s \\
Project demo & 90–180 s & 90–120 s \\
Deep breakdown & 3–8 min & separate video \\
Social teaser & 15–30 s & optional \\
Recruiter clip & 45–75 s & optional \\
\bottomrule
\end{longtable}
\end{center}




\subsection{Recommendation}




Create three video assets:




\begin{center}
\scriptsize
\begin{longtable}{>{\raggedright\arraybackslash}p{0.305\linewidth} >{\raggedright\arraybackslash}p{0.305\linewidth} >{\raggedright\arraybackslash}p{0.305\linewidth}}
\toprule
Asset & Duration & Use \\
\midrule
\endfirsthead
\toprule
Asset & Duration & Use \\
\midrule
\endhead
Main demo & 90–120 s & portfolio, LinkedIn Featured \\
Short teaser & 30–45 s & LinkedIn posts \\
Technical breakdown & 4–7 min & interviews, GitHub README \\
\bottomrule
\end{longtable}
\end{center}




\medskip\hrule\medskip




\subsection{7.2 Opening structure}




Weak opening:




\begin{mdcode}
Bloque de codigo: text
Logo fade-in → slow orbit → title after 20 seconds
\end{mdcode}




Strong opening:




\begin{mdcode}
Bloque de codigo: text
0–3 s: final polished viewer shot
3–6 s: value proposition
6–10 s: problem contrast
\end{mdcode}




Recommended opening caption:




\begin{mdcode}
Bloque de codigo: text
TwinSight X500
Unity WebGL technical visualization for drone assembly inspection
\end{mdcode}




Second caption:




\begin{mdcode}
Bloque de codigo: text
From static 2D documentation to interactive spatial understanding
\end{mdcode}




\medskip\hrule\medskip




\subsection{7.3 Shot pacing}




\begin{center}
\scriptsize
\begin{longtable}{>{\raggedright\arraybackslash}p{0.455\linewidth} >{\raggedright\arraybackslash}p{0.455\linewidth}}
\toprule
Shot type & Recommended length \\
\midrule
\endfirsthead
\toprule
Shot type & Recommended length \\
\midrule
\endhead
Hero reveal & 3–5 s \\
UI interaction & 4–6 s \\
Exploded view & 5–7 s \\
Clipping plane & 5–7 s \\
Visual modes grid & 5–8 s \\
CAD optimization comparison & 5–7 s \\
Metrics card & 4–6 s \\
End card & 3–5 s \\
\bottomrule
\end{longtable}
\end{center}




Avoid shots longer than 8 seconds unless the narration requires it.




\medskip\hrule\medskip




\section{8. TwinSight demo video architecture}




\subsection{8.1 Main 90-second version}




\begin{mdcode}
Bloque de codigo: text
00–05 s   Hook / hero shot
05–12 s   Problem: 2D documentation lacks spatial context
12–22 s   Solution: browser-based Unity WebGL viewer
22–34 s   Component selection and technical panels
34–46 s   Exploded view
46–58 s   Cross-section / clipping
58–68 s   Visual modes
68–78 s   CAD-to-realtime optimization
78–86 s   Evaluation metrics
86–90 s   Role, stack, links
\end{mdcode}




\subsection{8.2 120-second version}




\begin{mdcode}
Bloque de codigo: text
00–05 s    Hero shot
05–12 s    Project title and positioning
12–22 s    Problem context
22–32 s    WebGL browser interaction
32–44 s    Component selection
44–56 s    Exploded view
56–68 s    Cross-section
68–80 s    Visual modes
80–92 s    CAD optimization pipeline
92–102 s   UI Toolkit / runtime systems
102–112 s  Evaluation metrics
112–120 s  Links / stack / role
\end{mdcode}




\subsection{8.3 30-second teaser}




\begin{mdcode}
Bloque de codigo: text
00–03 s   Hero shot
03–07 s   Select component
07–12 s   Exploded view
12–17 s   Clipping
17–22 s   Visual modes
22–27 s   Metrics card
27–30 s   Portfolio link / title
\end{mdcode}




\medskip\hrule\medskip




\section{9. Required shot list for TwinSight}




\subsection{9.1 Core shots}




\begin{center}
\scriptsize
\begin{longtable}{>{\raggedright\arraybackslash}p{0.305\linewidth} >{\raggedright\arraybackslash}p{0.305\linewidth} >{\raggedright\arraybackslash}p{0.305\linewidth}}
\toprule
\# & Shot & Purpose \\
\midrule
\endfirsthead
\toprule
\# & Shot & Purpose \\
\midrule
\endhead
1 & Full drone hero view & first impression \\
2 & Browser frame / URL bar & WebGL proof \\
3 & Component hover/select & interaction proof \\
4 & Info panel opens & technical UI proof \\
5 & Exploded view activation & assembly understanding \\
6 & Exploded view close-up & hierarchy proof \\
7 & Clipping plane movement & inspection proof \\
8 & Internal structure revealed & technical visualization \\
9 & Visual modes grid & technical art proof \\
10 & X-ray / ghosted mode & inspection style \\
11 & Wireframe / blueprint mode & engineering aesthetic \\
12 & Thermal-style mode & visual experimentation \\
13 & CAD high-poly vs optimized & pipeline proof \\
14 & Triangle metric card & optimization proof \\
15 & SUS / NASA-TLX card & validation proof \\
16 & End card with stack & recruiter clarity \\
\bottomrule
\end{longtable}
\end{center}




\subsection{9.2 Optional shots}




\begin{center}
\scriptsize
\begin{longtable}{>{\raggedright\arraybackslash}p{0.455\linewidth} >{\raggedright\arraybackslash}p{0.455\linewidth}}
\toprule
Shot & Use \\
\midrule
\endfirsthead
\toprule
Shot & Use \\
\midrule
\endhead
Blender viewport & asset pipeline \\
Unity editor scene & implementation proof \\
UI Toolkit hierarchy & developer proof \\
Code snippet & C\# proof \\
mobile viewport & web/mobile proof \\
GitHub README & public proof \\
thesis evaluation sheet & academic validation \\
\bottomrule
\end{longtable}
\end{center}




\medskip\hrule\medskip




\section{10. Caption system}




\subsection{10.1 Caption style}




Use short captions, not paragraphs.




Good:




\begin{mdcode}
Bloque de codigo: text
Component selection + technical panel
\end{mdcode}




Bad:




\begin{mdcode}
Bloque de codigo: text
This system allows the user to select different drone components and view the associated information while navigating the interface.
\end{mdcode}




\subsection{10.2 Caption bank}




\begin{mdcode}
Bloque de codigo: text
Unity WebGL technical visualization
Browser-based drone assembly inspection
Component-level selection
Technical part information
Exploded assembly view
Cross-section inspection
Multiple visual modes
CAD-to-realtime optimization
95,617 optimized triangles
Evaluated with SUS + NASA-TLX
Built with Unity, C#, URP, UI Toolkit, Blender
\end{mdcode}




\subsection{10.3 Metrics captions}




\begin{mdcode}
Bloque de codigo: text
6.5M+ CAD triangle routes → 95,617 optimized triangles
12 validation participants
SUS average: 91.88
NASA-TLX Raw: 8.69 viewer vs 19.89 2D support
96 task-condition records
\end{mdcode}




Use “viewer vs 2D support” only if the methodology supports that comparison.




\medskip\hrule\medskip




\section{11. Voiceover recommendation}




\subsection{Main recommendation}




Use \textbf{captions + music only} for the public reel.




Add voiceover only for the long technical breakdown.




\subsection{Reason}




Recruiters often skim videos silently. Captions are safer.




\subsection{If using voiceover}




Keep it technical and restrained.




Example:




\begin{mdcode}
Bloque de codigo: text
TwinSight X500 is a Unity WebGL technical visualization prototype for drone assembly inspection. It converts CAD-derived assets into an optimized browser-based viewer with component selection, exploded view, cross-section tools, technical panels, and validation through SUS and NASA-TLX.
\end{mdcode}




\medskip\hrule\medskip




\section{12. Technical breakdown video}




Create a separate video, not the main reel.




\subsection{Recommended length}




\begin{mdcode}
Bloque de codigo: text
4–7 minutes
\end{mdcode}




\subsection{Structure}




\begin{mdcode}
Bloque de codigo: text
1. Problem and thesis context
2. CAD source and optimization constraints
3. Blender cleanup / retopology / baking
4. Unity WebGL architecture
5. Runtime interaction systems
6. UI Toolkit interface
7. Visual modes
8. Evaluation methodology
9. Limitations
10. Next steps
\end{mdcode}




\subsection{What to show}




\begin{itemize}
\item screen capture from Unity editor;
\item scene hierarchy;
\item C\# scripts, briefly;
\item shader graph, briefly;
\item profiler if available;
\item before/after asset counts;
\item GitHub README;
\item final browser build.
\end{itemize}




\medskip\hrule\medskip




\section{13. Editing rules}




\subsection{13.1 Must do}




\begin{itemize}
\item show final output in first 5 seconds;
\item use captions on every feature;
\item keep shots short;
\item show interaction, not just orbit shots;
\item include metrics;
\item include stack;
\item include one clear end card;
\item export cleanly at 1080p minimum;
\item use consistent typography;
\item avoid excessive transitions.
\end{itemize}




\subsection{13.2 Must avoid}




\begin{itemize}
\item long intro logos;
\item slow cinematic-only orbit;
\item unclear UI text;
\item music louder than explanation;
\item too many effects;
\item pretending it is a shipped commercial product;
\item hiding that it is an academic/portfolio prototype;
\item overusing “AI” in the main reel;
\item adding unimplemented XR/IoT/WebAR features.
\end{itemize}




\medskip\hrule\medskip




\section{14. Export specifications}




\subsection{Main reel}




\begin{center}
\scriptsize
\begin{longtable}{>{\raggedright\arraybackslash}p{0.455\linewidth} >{\raggedright\arraybackslash}p{0.455\linewidth}}
\toprule
Parameter & Recommendation \\
\midrule
\endfirsthead
\toprule
Parameter & Recommendation \\
\midrule
\endhead
Resolution & 1920×1080 \\
Frame rate & 30 fps or 60 fps \\
Duration & 90–120 s \\
Format & MP4 H.264 \\
Audio & optional music \\
Captions & burned-in \\
Aspect & 16:9 \\
Primary use & portfolio, LinkedIn, README \\
\bottomrule
\end{longtable}
\end{center}




\subsection{LinkedIn teaser}




\begin{center}
\scriptsize
\begin{longtable}{>{\raggedright\arraybackslash}p{0.455\linewidth} >{\raggedright\arraybackslash}p{0.455\linewidth}}
\toprule
Parameter & Recommendation \\
\midrule
\endfirsthead
\toprule
Parameter & Recommendation \\
\midrule
\endhead
Resolution & 1080×1080 or 1920×1080 \\
Duration & 30–45 s \\
Captions & mandatory \\
Hook & first 2 seconds \\
CTA & portfolio link / project name \\
\bottomrule
\end{longtable}
\end{center}




\subsection{Vertical teaser}




\begin{center}
\scriptsize
\begin{longtable}{>{\raggedright\arraybackslash}p{0.455\linewidth} >{\raggedright\arraybackslash}p{0.455\linewidth}}
\toprule
Parameter & Recommendation \\
\midrule
\endfirsthead
\toprule
Parameter & Recommendation \\
\midrule
\endhead
Resolution & 1080×1920 \\
Duration & 20–30 s \\
Use & optional \\
Risk & UI may be hard to read \\
\bottomrule
\end{longtable}
\end{center}




\medskip\hrule\medskip




\section{15. Recommended demo reel script}




\subsection{15.1 90-second caption script}




\begin{mdcode}
Bloque de codigo: text
[00–05]
TwinSight X500
Unity WebGL technical visualization for drone assembly inspection

[05–12]
Problem:
2D documentation forces users to mentally reconstruct spatial relationships.

[12–22]
Solution:
An interactive browser-based 3D viewer for technical inspection.

[22–34]
Select components.
Highlight relationships.
Open technical part information.

[34–46]
Exploded view:
understand assembly hierarchy and spatial placement.

[46–58]
Cross-section tools:
inspect hidden structures and internal relationships.

[58–68]
Visual modes:
realistic, X-ray, ghosted, blueprint, wireframe, thermal-style.

[68–78]
CAD-to-realtime pipeline:
optimized complex CAD-derived assets for WebGL.

[78–86]
Evaluation:
SUS 91.88
NASA-TLX Raw: 8.69 viewer vs 19.89 2D support
12 participants

[86–90]
Unity · C# · WebGL · URP · UI Toolkit · Blender
Built by Alexander Woodcock Salomón
\end{mdcode}




\subsection{15.2 Shorter LinkedIn teaser script}




\begin{mdcode}
Bloque de codigo: text
Static 2D assembly docs are hard to understand.

TwinSight X500 turns drone assembly into an interactive Unity WebGL viewer.

Select components.
Explode the assembly.
Cut through structures.
Switch visual modes.
Inspect technical data.

CAD-derived assets optimized for browser-based real-time 3D.
\end{mdcode}




\medskip\hrule\medskip




\section{16. Portfolio embedding strategy}




\subsection{Where to place the video}




\begin{center}
\scriptsize
\begin{longtable}{>{\raggedright\arraybackslash}p{0.455\linewidth} >{\raggedright\arraybackslash}p{0.455\linewidth}}
\toprule
Location & Video \\
\midrule
\endfirsthead
\toprule
Location & Video \\
\midrule
\endhead
Portfolio homepage & 30–45 s teaser \\
TwinSight case study top & 90–120 s main demo \\
GitHub README & compressed GIF + YouTube/Vimeo link \\
LinkedIn Featured & main video or case-study link \\
ArtStation & visual breakdown and short video \\
Interview follow-up & technical breakdown \\
\bottomrule
\end{longtable}
\end{center}




\subsection{Recommended hosting}




\begin{center}
\scriptsize
\begin{longtable}{>{\raggedright\arraybackslash}p{0.455\linewidth} >{\raggedright\arraybackslash}p{0.455\linewidth}}
\toprule
Platform & Use \\
\midrule
\endfirsthead
\toprule
Platform & Use \\
\midrule
\endhead
YouTube unlisted/public & main video \\
Vimeo & polished portfolio embed \\
GitHub README & GIF/video thumbnail \\
LinkedIn native upload & short teaser \\
Portfolio site & embedded main reel \\
\bottomrule
\end{longtable}
\end{center}




\medskip\hrule\medskip




\section{17. GitHub README video integration}




Use this structure:




\begin{mdcode}
Bloque de codigo: markdown
## Demo

[![TwinSight X500 Demo](./media/twinsight-demo-thumbnail.jpg)](https://youtube.com/...)

## Highlights

- Unity WebGL technical visualization
- CAD-to-realtime optimization
- Component selection and technical UI
- Exploded view and clipping tools
- Visual modes for inspection
- SUS + NASA-TLX evaluation
\end{mdcode}




Add a lightweight GIF only if file size is controlled.




\medskip\hrule\medskip




\section{18. ArtStation video integration}




ArtStation should not be only a final render gallery.




Use:




\begin{mdcode}
Bloque de codigo: text
1. 30–60 s video
2. final beauty screenshots
3. visual modes grid
4. topology / optimization screenshots
5. UI screenshots
6. breakdown captions
7. tools used
8. role and contribution
\end{mdcode}




ArtStation tags:




\begin{mdcode}
Bloque de codigo: text
Unity
Technical Art
Realtime
WebGL
Blender
CAD Optimization
Visualization
Drone
Interactive 3D
\end{mdcode}




\medskip\hrule\medskip




\section{19. What this changes in the roadmap}




\subsection{Keep}




\begin{mdcode}
Bloque de codigo: text
TwinSight remains the flagship.
\end{mdcode}




\subsection{Add}




\begin{mdcode}
Bloque de codigo: text
TwinSight needs a polished video asset before heavy applications.
\end{mdcode}




\subsection{Priority}




The video is not optional. It is one of the highest-leverage artifacts because:




\begin{itemize}
\item recruiters skim;
\item technical managers need proof fast;
\item Unity/WebGL projects are hard to understand from text only;
\item the candidate lacks formal industry experience;
\item the video can compensate by showing proof.
\end{itemize}




\medskip\hrule\medskip




\section{20. Minimum viable demo video checklist}




Before applications, produce at least:




\begin{mdcode}
Bloque de codigo: text
[ ] 30–45 s teaser
[ ] 90–120 s main demo
[ ] 10–15 strong screenshots
[ ] GitHub README embed
[ ] portfolio case-study embed
[ ] LinkedIn Featured entry
\end{mdcode}




\subsection{Acceptance criteria}




The main demo is acceptable if a viewer understands within 20 seconds:




\begin{mdcode}
Bloque de codigo: text
what it is,
what problem it solves,
what the user can do,
why it is technically relevant.
\end{mdcode}




\medskip\hrule\medskip




\section{21. Next recommended task}




After this module, proceed with:




\begin{mdcode}
Bloque de codigo: text
F3_artstation_breakdown_benchmark
\end{mdcode}




\subsection{Tool recommendation}




\begin{center}
\scriptsize
\begin{longtable}{>{\raggedright\arraybackslash}p{0.455\linewidth} >{\raggedright\arraybackslash}p{0.455\linewidth}}
\toprule
Tool & Use \\
\midrule
\endfirsthead
\toprule
Tool & Use \\
\midrule
\endhead
Agent mode & Yes \\
Deep Research & No \\
\bottomrule
\end{longtable}
\end{center}




\subsection{Reason}




The demo reel benchmark defines video structure. The next missing piece is the \textbf{visual breakdown style} for ArtStation / portfolio pages, especially:




\begin{itemize}
\item topology;
\item materials;
\item optimization;
\item Blender portrait;
\item technical art breakdown;
\item visual modes;
\item real-time asset presentation.
\end{itemize}




\subsection{Prompt for next task}




\begin{mdcode}
Bloque de codigo: text
Execute F3_artstation_breakdown_benchmark.

Objective:
Collect ArtStation or portfolio breakdowns useful for Alexander’s Blender portrait, TwinSight X500 and technical art presentation.

Find 15–25 examples involving:
- topology breakdown;
- material breakdown;
- grooming breakdown;
- lighting breakdown;
- shader breakdown;
- real-time asset breakdown;
- before/after optimization;
- character study presentation;
- technical art process documentation;
- Unity technical art breakdown;
- real-time visualization breakdown.

Output structured tables only.

Fields:
- Project label
- URL
- Platform
- Breakdown type
- Tools shown
- Images used
- Video/GIF used
- Technical text depth
- Tags
- What Alexander can adapt
- What Alexander should not copy
- Confidence
- Date checked

Rules:
- Public pages only.
- No invented data.
- If unavailable, write “not found”.
- Keep cells short.
- Tables only.
- Include URLs.
- No strategy.
\end{mdcode}



